\section*{Capítulo \ref{ch:g9:periodic-table-first-look} --- La tabla periódica: un primer vistazo}

\begin{solution}{exo:g9:periodic-table-first-look:1}
Por orden creciente de \omterm{def:g9:inside-the-atom:atomic-number}{número atómico} $Z$.
\end{solution}

\begin{solution}{exo:g9:periodic-table-first-look:2}
Carbono: \omterm{def:g9:periodic-table-first-look:period}{periodo} 2, \omterm{def:g9:periodic-table-first-look:period}{grupo} 14. Magnesio: \omterm{def:g9:periodic-table-first-look:period}{periodo} 3, \omterm{def:g9:periodic-table-first-look:period}{grupo} 2. Argón:
\omterm{def:g9:periodic-table-first-look:period}{periodo} 3, \omterm{def:g9:periodic-table-first-look:period}{grupo} 18.
\end{solution}

\begin{solution}{exo:g9:periodic-table-first-look:3}
\omterm{def:g9:periodic-table-first-look:metal}{Metales}: el hierro, el cobre, el aluminio. \omterm{def:g9:periodic-table-first-look:metal}{No metales}: el azufre, el
oxígeno, el cloro.
\end{solution}

\begin{solution}{exo:g9:periodic-table-first-look:4}
El potasio (\ce{K}) debajo del sodio; el bromo (\ce{Br}) debajo del
cloro.
\end{solution}

\begin{solution}{exo:g9:periodic-table-first-look:5}
El calcio (\ce{Ca}, $Z = 20$) y el nitrógeno (\ce{N}, $Z = 7$).
\end{solution}

\begin{solution}{exo:g9:periodic-table-first-look:6}
Está en el mismo grupo que el sodio, el \omterm{def:g9:periodic-table-first-look:period}{grupo} 1, y los elementos de un
grupo se comportan igual.
\end{solution}

\begin{solution}{exo:g9:periodic-table-first-look:7}
Un \omterm{def:g9:periodic-table-first-look:metal}{metal}: es brillante, conduce la electricidad y forma un \omterm{def:g9:ions:ion}{ion} positivo.
\end{solution}

\begin{solution}{exo:g9:periodic-table-first-look:8}
Para mantener en la misma columna los elementos de propiedades
parecidas, aunque eso obligara a dejar una casilla vacía. Los huecos
eran predicciones: los elementos que faltaban se encontraron después con
las propiedades que él había descrito.
\end{solution}

\begin{solution}{exo:g9:periodic-table-first-look:9}
Carbono: $\ce{C} = 12$; oxígeno: $\ce{O} = 16$. Los huecos ``? $= 68$''
(junto a Al $= 27.4$) y ``? $= 70$'' (junto a Si $= 28$).
\end{solution}

\begin{solution}{exo:g9:periodic-table-first-look:10}
\omterm{def:g9:periodic-table-first-look:period}{Periodo} 2: 8 elementos (del litio al neón). \omterm{def:g9:periodic-table-first-look:period}{Periodo} 4: 18 elementos (del
potasio al criptón).
\end{solution}

\begin{solution}{exo:g9:periodic-table-first-look:11}
\omterm{def:g9:periodic-table-first-look:period}{Grupo} 18. Allí donde nada debe reaccionar: el argón de las bombillas y
de los gases de soldadura, el helio de los globos (no \omterm{def:g4:what-a-fire-needs:burning}{arde}).
\end{solution}

\begin{solution}{exo:g9:periodic-table-first-look:12}
$(27.0 + 114.8)/2 = 70.9$, cerca de 69.7.
\end{solution}

\begin{solution}{pb:g9:periodic-table-first-look:1}
\textbf{1.} $Z = 32$; \omterm{def:g9:periodic-table-first-look:period}{periodo} 4, \omterm{def:g9:periodic-table-first-look:period}{grupo} 14.

\textbf{2.} Encima: el silicio. Debajo: el estaño. A la izquierda: el
galio. A la derecha: el arsénico.

\textbf{3.} Un metaloide.

\textbf{4.} Está en el mismo grupo que el silicio, y los elementos de un
grupo se comportan igual.

\textbf{5.} $(28.1 + 118.7)/2 = 73.4$.

\textbf{6.} $(69.7 + 74.9)/2 = 72.3$.

\textbf{7.} La media izquierda--derecha (72.3): a lo largo de un \omterm{def:g9:periodic-table-first-look:period}{periodo}
el peso atómico crece de forma regular, elemento a elemento, mientras
que al bajar por un grupo da grandes saltos.

\textbf{8.} $72.6 - 70 = 2.6$.

\textbf{9.} $(28.1 + 118.7 + 69.7 + 74.9)/4 = 291.4/4 = 72.85$, es
decir, 72.9.

\textbf{10.} La media, 72.9, está cerca tanto del 72 de Mendeléiev como
del 72.3 de Winkler (y del 72.6 actual).

\textbf{11.} Un elemento predicho antes de que nadie lo hubiera visto se
encontró con las propiedades anunciadas: la tabla no era solo una lista,
podía hacer predicciones.

\textbf{12.} Alrededor de 72.9.
\end{solution}
